\documentclass[11pt]{amsart}

\usepackage[T1]{fontenc}
\usepackage{lmodern}
\usepackage{microtype}
\usepackage{amsmath,amssymb,amsthm}
\usepackage{booktabs}
\usepackage[hidelinks]{hyperref}

\hypersetup{
  pdftitle={Positive Generalized Corner-Vector Designs on S3 Have Degree at Most 13},
  pdfsubject={Spherical designs and positive cubature},
  pdfkeywords={spherical design, cubature formula, hyperoctahedral group, generalized corner vector, Bernstein basis}
}

\newtheorem{theorem}{Theorem}
\newtheorem{lemma}[theorem]{Lemma}
\newtheorem{corollary}[theorem]{Corollary}
\theoremstyle{remark}
\newtheorem{remark}[theorem]{Remark}

\newcommand{\Sph}{\mathbb{S}}
\newcommand{\mean}[1]{\left\langle #1\right\rangle}

\title[Generalized Corner-Vector Designs on $\Sph^3$]
{Positive Generalized Corner-Vector Designs on
\texorpdfstring{$\Sph^3$}{S3} Have Degree at Most 13}

\date{}

\subjclass[2020]{Primary 65D32; Secondary 05B30}
\keywords{spherical design, cubature formula, hyperoctahedral group,
generalized corner vector, separating polynomial, Bernstein basis}

\begin{document}

\begin{abstract}
Tanino, Tamaru, Hirao, and Sawa have proved that positive weighted
spherical designs obtained from generalized corner-vector orbits in
dimension four have degree at most $15$, and asked whether a weighted
$15$-design on $\Sph^3$ with more than five proper orbits exists. We
give a $B_4$-invariant polynomial of degree $14$ that is positive on
every generalized corner-vector orbit but has negative spherical mean.
Consequently, every positive design of this type on $\Sph^3$ has degree
at most $13$. This gives a negative answer, for positive weights, to
their Problem~6.1, independently of the number of proper orbits.
\end{abstract}

\maketitle

\section{The obstruction}

Let $B_4$ act on $\mathbb{R}^4$ by signed coordinate permutations. For
$a>0$ and $s\in\{0,1,2,3\}$, set
\[
v_{a,s}
=
\frac{1}{\sqrt{a^2+s}}
\bigl(a,\underbrace{1,\ldots,1}_{s},0,\ldots,0\bigr)
\in\Sph^3,
\]
and denote its $B_4$-orbit by $v_{a,s}^{B_4}$. These are the
generalized corner-vector orbits of Tanino--Tamaru--Hirao--Sawa
\cite[Definition~2.19]{TTHS}; the orbit is called proper when $a\ne1$.
They proved that a positive weighted design assembled from such orbits
has degree at most $15$ in dimension four \cite[Theorem~3.1]{TTHS}, and
they asked \cite[Problem~6.1]{TTHS}:
\begin{quote}
Does there exist a weighted $15$-design on $\Sph^3$ of type~(1) with
more than $5$ proper orbits?
\end{quote}
Their ``type~(1)'' is the display numbered~(1) on p.~1388 of
\cite{TTHS}: a cubature formula for normalized surface measure whose
node set is a union of generalized corner-vector orbits, carrying one
weight per orbit. That is the class of formula~\eqref{eq:cubature}
below, and it is also the class occurring in the hypothesis of
\cite[Theorem~3.1]{TTHS}, whose formula~(11) has the same shape. Let
$\sigma$ denote normalized surface measure on $\Sph^3$.

\begin{theorem}\label{thm:main}
Let $m\ge1$, $a_i>0$, $s_i\in\{0,1,2,3\}$, and $W_i>0$. The formula
\begin{equation}\label{eq:cubature}
\int_{\Sph^3} f\,d\sigma
=
\sum_{i=1}^{m}W_i\sum_{x\in v_{a_i,s_i}^{B_4}}f(x)
\end{equation}
cannot be exact for every polynomial $f$ of degree at most $14$.
\end{theorem}

\begin{corollary}\label{cor:problem}
Every positive weighted spherical design of generalized corner-vector
type on $\Sph^3$ has degree at most $13$. In particular,
Problem~6.1 of \cite{TTHS} has a negative answer for positive weights,
as required of a weighted design by \cite[Definition~2.2]{TTHS},
regardless of the number of proper orbits.
\end{corollary}

Put $y_i=x_i^2$, and let $e_j=e_j(y_1,y_2,y_3,y_4)$ be the $j$th
elementary symmetric polynomial. Define
\begin{align}\label{eq:p}
p={}&
2+240e_2+1345e_3+227924e_4
-3518e_2^2-448e_2e_3+283604e_3^2 \notag\\
&-1529173e_2e_4+12886e_2^3
+4855151e_3e_4-132972e_2^2e_3.
\end{align}
Then $p$ is $B_4$-invariant and has degree $14$ in
$x_1,\ldots,x_4$. For a polynomial $f$ on $\Sph^3$, write
\[
\mean{f}=\int_{\Sph^3}f\,d\sigma.
\]

\begin{lemma}\label{lem:mean}
The spherical mean of $p$ is negative:
\[
\mean{p}=-\frac{60763}{143360}<0.
\]
\end{lemma}

\begin{proof}
For a uniformly distributed point of $\Sph^3$, the vector
$(y_1,y_2,y_3,y_4)$ has the Dirichlet distribution with parameters
$(1/2,1/2,1/2,1/2)$. Hence, for
$\alpha_i\in\mathbb{Z}_{\ge0}$,
\[
\mean{y_1^{\alpha_1}\cdots y_4^{\alpha_4}}
=
\frac{\prod_{i=1}^{4}(1/2)_{\alpha_i}}
     {(2)_{\alpha_1+\cdots+\alpha_4}},
\qquad (z)_0=1.
\]
Expanding \eqref{eq:p} and applying this identity gives the stated
value.
\end{proof}

\begin{lemma}\label{lem:positive}
For every $a>0$ and $s\in\{0,1,2,3\}$,
\[
p(v_{a,s})>0.
\]
\end{lemma}

\begin{proof}
For $s=0$, one has $v_{a,0}=(1,0,0,0)$ and $p(v_{a,0})=2$.
Let $s\in\{1,2,3\}$ and set
\[
u=\frac{a^2}{a^2+s}\in(0,1),
\qquad
c=\frac{1-u}{s}.
\]
Since $p$ is a symmetric function of the $y_i$, we may take the squared
coordinates of $v_{a,s}$ to be
$(u,c,\ldots,c,0,\ldots,0)$, with $c$ repeated $s$ times. Direct
substitution into \eqref{eq:p} gives
\[
p(v_{a,1})=H_1(u),\qquad
32p(v_{a,2})=H_2(u),\qquad
729p(v_{a,3})=H_3(u),
\]
where
\begin{align*}
H_1(u)={}&
-12886u^6+38658u^5-42176u^4+19922u^3
-3758u^2+240u+2,\\
H_2(u)={}&
-598374u^7+2387827u^6-3979288u^5+3540713u^4\\
&-1750294u^3+440357u^2-42268u+1391,\\
H_3(u)={}&
-51606520u^7+213893327u^6-355558770u^5+299081978u^4\\
&-130946506u^3+26917938u^2-1819702u+39713.
\end{align*}

We certify positivity by exact Bernstein coefficients. For an interval
$I=[r,q]\subseteq[0,1]$ and a polynomial $H$ of degree $d$, write
\[
H(r+(q-r)z)
=
\sum_{k=0}^{d}\beta_{I,k}(H)
\binom{d}{k}z^k(1-z)^{d-k}.
\]
If $H(r+(q-r)z)=\sum_{j=0}^{d}c_jz^j$, then
\begin{equation}\label{eq:bernstein}
\beta_{I,k}(H)
=
\sum_{j=0}^{k}c_j
\frac{\binom{k}{j}}{\binom{d}{j}}.
\end{equation}
For each $s$, the following table gives a subdivision
$0=r_0<r_1<r_2<r_3=1$ and the minimum Bernstein coefficient on each
successive interval, where $d=\deg H_s$, so that $d=6$ for $s=1$ and
$d=7$ for $s=2,3$. The coefficients \eqref{eq:bernstein} depend on $d$;
in particular $H_1$ is not to be padded to degree $7$.

\begin{center}
\small
\renewcommand{\arraystretch}{1.35}
\begin{tabular}{@{}c c c@{}}
\toprule
$s$ & $(r_0,r_1,r_2,r_3)$
& $\bigl(\min_k\beta_{[r_{j-1},r_j],k}(H_s)\bigr)_{j=1}^{3}$ \\
\midrule
$1$ & $\left(0,\frac16,\frac56,1\right)$
& $\left(\frac{14287}{7776},\frac{46531}{23328},
          \frac{14287}{7776}\right)$ \\[1.5ex]
$2$ & $\left(0,\frac1{12},\frac34,1\right)$
& $\left(\frac{137627845}{2985984},
          \frac{414113333}{5971968},\frac{219941}{5376}\right)$ \\[1.5ex]
$3$ & $\left(0,\frac1{20},\frac35,1\right)$
& $\left(\frac{542097724127}{448000000},
          \frac{17317137}{15625},\frac{4479262}{3125}\right)$ \\
\bottomrule
\end{tabular}
\end{center}

Every displayed number is positive. Since the Bernstein basis
functions are nonnegative on $[0,1]$ and sum to $1$, it follows that
$H_s(u)>0$ on every listed interval, hence on $[0,1]$. Therefore
$p(v_{a,s})>0$ for $s=1,2,3$ as well.
\end{proof}

\begin{proof}[Proof of Theorem~\ref{thm:main}]
Assume that \eqref{eq:cubature} is exact through degree $14$. Applying
it to the polynomial $p$ and using its $B_4$-invariance gives
\[
-\frac{60763}{143360}
=
\mean{p}
=
\sum_{i=1}^{m}
W_i\lvert v_{a_i,s_i}^{B_4}\rvert p(v_{a_i,s_i}).
\]
The right-hand side is strictly positive by $W_i>0$ and
Lemma~\ref{lem:positive}, contradicting Lemma~\ref{lem:mean}.
\end{proof}

\begin{remark}
Since $-I\in B_4$, every formula of the form \eqref{eq:cubature} is
antipodal. Exactness through degree $14$ is therefore equivalent to
exactness through degree $15$.

Positivity is essential, and the positive-weight reading is the one
intended in \cite{TTHS}: a weighted design there carries by definition a
weight function with values in $\mathbb{R}_{>0}$
\cite[Definition~2.2]{TTHS}, so the hypothesis $W_i>0$ of
Theorem~\ref{thm:main} is theirs and not a restriction added here. The
displays \cite[(1) and~(11)]{TTHS} do not repeat the sign condition; that
the positive reading is nonetheless the intended one is confirmed by
\cite[\S6]{TTHS}, where it is recorded, immediately before Problem~6.1
is posed, that $15$-designs on $\Sph^3$ with negative weights do exist.
Theorem~\ref{thm:main} accordingly says nothing about signed cubature
formulas, nor about node families outside the generalized corner-vector
construction. Degree $13$ is attained within this class: \cite[\S6]{TTHS}
reports two weighted $13$-designs on $\Sph^3$ of this type, each with
five orbits and positive weights, with nodes and weights given there
numerically to six significant figures. On the evidence of those
examples the bound of Corollary~\ref{cor:problem} is sharp; verifying
them exactly is not attempted here.
\end{remark}

\paragraph{Exact verification.}
Everything above is finite and exact, and is recomputable from the data
printed here; the only mathematical input is the eleven coefficients
of \eqref{eq:p}. From them one expands $p$ into its $130$ monomials in
the $y_i$ and confirms its $B_4$-invariance and its degree $14$ in
$x_1,\ldots,x_4$; evaluates $\mean{p}=-60763/143360$ by two routes,
direct expansion against the Dirichlet moment functional of
Lemma~\ref{lem:mean} and a harmonic route through the Laplacian, the
moment functional itself being cross-validated against the $24$-cell
through degree $5$ and against an iterated-Laplacian formula on all
$330$ even monomials of degree at most $14$; substitutes the profiles
of Lemma~\ref{lem:positive} to recover $H_1,H_2,H_3$ with their
integrality scalings $1$, $32$, $729$; verifies the identity
\eqref{eq:bernstein} itself, rather than assuming it, and with it all
nine minimum Bernstein coefficients of the table; and re-certifies
positivity of $H_1,H_2,H_3$ on $[0,1]$ a second way, by exact Sturm
root counting. No floating-point arithmetic enters at any stage. This
computation was re-run independently in exact rational arithmetic, and
all $48$ of its checks agree with the numbers printed here. The program
and its output log accompany this note; nothing above depends on them,
since the Bernstein table is itself the certificate.

\medskip

\paragraph{Scope.}
Two limits belong with that statement. First, the claim quantifies over
the continuum $a>0$, so no census is possible and none was attempted:
the universal statement is discharged exactly by the Bernstein
certificate over the tabulated partition, re-certified by Sturm, with
exact evaluation of $p$ at $609$ genuine orbit points as a further
check. Second, agreement is what is claimed for the independent re-run,
not a second proof: a few of its checks are self-tests of its own
construction, which no object could fail, and the agreement of the
Pochhammer and double-factorial forms of the Dirichlet moment is an
algebraic rewriting rather than an independent derivation.

\begin{thebibliography}{1}

\bibitem{TTHS}
K.~Tanino, T.~Tamaru, M.~Hirao, and M.~Sawa,
\newblock More on the corner-vector construction for spherical designs,
\newblock \emph{Algebraic Combinatorics} \textbf{8} (2025), no.~5,
1387--1414,
\newblock \href{https://doi.org/10.5802/alco.441}{doi:10.5802/alco.441}.

\end{thebibliography}

\end{document}